\label{sec:v71-lens-comparison}
A second comparison is with smooth obstacle scattering.  In the corrected
planar theorem of Noakes--Stoyanov \cite[Theorem~1.1]{NoakesStoyanov2020},
a finite disjoint union of strictly convex compact domains with $C^3$
boundaries is determined globally by almost-everywhere travelling times,
or by almost-everywhere scattering length spectra.  For the former data,
one fixes an enclosing circle and records the exit time for each incoming
position and direction.  For the latter, the data are sets of sojourn times
indexed by the incoming and outgoing directions.  These are labelled
families, not an unlabelled list of lengths.  The cited paper supplies a
separate two-dimensional proof; its earlier higher-dimensional argument
alone does not cover the planar case.  Smooth recovery from rich scattering
data is therefore not a new category of rigidity supplied by this paper.

The distinctions relevant to Theorem~\ref{thm:v70-main} are concrete.
The exterior geometry there is a finite compact obstacle, whereas here it
is a periodic table with a selected recurrent itinerary.  An enclosing
Euclidean observation frame in the scattering problem does not supply
internal reflecting contacts.  Our vertices, tangent frames, normals and
phase labels are supplied, while the reflecting graphs remain unknown.
Our data are two conditional endpoint distributions at each marked phase,
not exit times indexed by all exterior incoming rays.  Spatially, our
smooth conclusion is the collection of contact germs, not the boundaries
away from the visited collars.  The smoothness and quantitative priors
required for our finite experiment are also stronger than the $C^3$
uniqueness hypothesis just cited.  Thus neither regularity nor spatial
extent provides a dominance claim over that theorem.

There is an exact local reason for these differences, not merely a
terminological one.  Proposition~\ref{prop:v71-locality} proves that all
finite-flight laws of the specified experiment are unchanged when the
unvisited boundary is altered while the local flight geometry is fixed.
With free area fixed, this includes the success probabilities and the
failure atom, not only the conditional densities.  Conversely the two
offset limits determine the complete contact germs by the smooth inverse.
Corollary~\ref{cor:v71-local-completion} realizes both assertions in
nonconstant equal-area families of closed obstacles.  On a common
compact-obstacle realization, their exterior scattering data distinguish
the different completions.  Equality of our local observation therefore
does not supply the global-data equality hypothesis of the planar lens
theorem.  No assertion in the reverse direction, or of two-way
incomparability, is needed.

The contribution is consequently a local sufficiency statement at a
singular observation scale.  The cofactor identity alone describes a
linear reference; the two-ended determinant estimate must show that the
\emph{nonlinear physical numerator} remains relatively controlled after
its exponentially small reference has been removed.  Once the resulting
laws yield actual half-line actions, coefficient inversion still leaves
every flat perturbation unresolved.  The summable stationary envelope and
constant-one contraction of actual visits supply the functional step that
identifies that remaining information.  The finite experiment reaches the
finite differentiated topology used by this step while charging the rare
successes.  This chain, rather than smooth rigidity in isolation, is the
mathematical gain.  Its ingredients and observation hypotheses are not
interchangeable with a supplied global scattering relation.  The locality
comparison identifies the object determined by the chain; it is not an
additional claim of global reconstruction or an independent statistical
method.
